% !TEX program = xelatex
\documentclass[11pt]{article}
\usepackage{fontspec}
\usepackage{xeCJK}
\setmainfont{CMU Serif}
\newcommand{\setcjkserifandmono}[1]{%
  \setCJKmainfont{#1}%
  \setCJKmonofont{#1}%
}
\IfFontExistsTF{Noto Serif CJK SC}{\setcjkserifandmono{Noto Serif CJK SC}}{%
  \IfFontExistsTF{Source Han Serif SC}{\setcjkserifandmono{Source Han Serif SC}}{%
    \setcjkserifandmono{FandolSong-Regular}%
  }%
}%
\usepackage[margin=1in]{geometry}
\usepackage{amsmath}
\usepackage{amssymb}
\usepackage{booktabs}
\usepackage{array}
\usepackage{enumitem}
\usepackage{xcolor}
\usepackage{hyperref}
\hypersetup{colorlinks=true, linkcolor=blue, urlcolor=blue}

\setlength{\parindent}{0pt}
\setlength{\parskip}{0.45\baselineskip}
\setlength{\emergencystretch}{2em}
\setlength{\tabcolsep}{4pt}
\renewcommand{\arraystretch}{1.08}
\setlist[itemize]{leftmargin=1.6em, itemsep=2pt, topsep=3pt}

\newcommand{\sid}{\mathrm{SID}}
\newcommand{\rqvae}{\mathrm{RQ\mbox{-}VAE}}
\newcommand{\sg}{\mathrm{sg}}
\newcommand{\row}{\mathrm{RowNorm}}
\newcommand{\cosim}{\mathrm{cos}}
\newcommand{\pos}{\mathcal{P}}
\newcommand{\negset}{\mathcal{N}}
\newcommand{\lmh}{\mathrm{LMH}}

\begin{document}

{\LARGE SID 多层协同信息注入当前实验汇报}

\vspace{0.4\baselineskip}

\section{摘要}

我们当前研究的核心问题是：\textbf{如何把协同层级信息融入 SID tokenizer 的构建，使下游大语言模型更容易学习到正确的推荐范式？}

目前的成果是：我们在 RQ-VAE tokenizer 中对 L1/L2/L3 分别施加不同粒度的约束：L1 保持语义近邻路由，L2 注入多跳协同邻域信号，L3 保持局部转移精修。该实验在 SFT 阶段，NDCG 在 @1/@3/@5/@10 四个截止点上全面超过 Baseline SFT（基线监督微调）：

\[
\begin{array}{c|cccc}
& @1 & @3 & @5 & @10 \\ \hline
\text{Baseline} & 0.06706 & 0.08501 & 0.09315 & 0.10372 \\
\text{本文方法} & 0.07059 & 0.08813 & 0.09489 & 0.10438 \\
\end{array}
\]

头部提升明显（@1 +5.3\%），说明协同层级注入已经开始改善下游排序质量。

\section{Motivation：为什么 SID 需要协同层级信息}

MiniOneRec 的基本思路是先用 item semantic embedding（物品语义嵌入）训练 RQ-VAE tokenizer，然后把每个物品表示成三层 SID：

\[
z_i = \bigl(z_i^{(1)}, z_i^{(2)}, z_i^{(3)}\bigr),
\]

下游 LLM（大语言模型）根据用户历史自回归预测目标物品的 SID。原版 tokenizer 的输入主要来自 title/description embedding（标题/描述嵌入），因此 SID 空间天然偏语义。

问题在于：\textbf{语义相近不一定协同相近，协同相近也不一定语义相近。} 例如两个工业耗材在标题上非常相似，但用户购买场景可能完全不同；反过来，一个工具和一种配件在语义文本上差异明显，却可能经常出现在同一类购买序列中。纯语义 SID 会让下游模型在某些位置学习到不合适的 item family（物品族）或 route（路由）。

因此，我们的研究目标是构造一个更适合下游学习的层级 SID 空间：

\begin{itemize}
  \item L1 应该保持语义近邻路由，提供稳定的第一层分区。
  \item L2 应该接收多跳协同邻域信息，在语义分区内做协同重排。
  \item L3 应该接收 局部转移精修信息，减少高相似度物品混淆。
\end{itemize}

\section{原始 RQ-VAE SID 构建}

给定物品语义嵌入 $x_i \in \mathbb{R}^d$，RQ-VAE 先编码得到连续表示，然后通过三层 residual quantization（残差量化）得到量化向量 $q_i^{(1)},q_i^{(2)},q_i^{(3)}$。原版重建表示为：

\[
h_i = q_i^{(1)} + q_i^{(2)} + q_i^{(3)}, \qquad \hat{x}_i = f_{\mathrm{dec}}(h_i).
\]

基础训练目标可以写成：

\[
\mathcal{L}_{\mathrm{RQ}} =
\mathcal{L}_{\mathrm{rec}}
+ \lambda_q \mathcal{L}_{\mathrm{quant}},
\]

其中

\[
\mathcal{L}_{\mathrm{rec}}
= \frac{1}{N}\sum_{i=1}^{N}\lVert \hat{x}_i-x_i\rVert_2^2.
\]

该目标能较好保持 semantic reconstruction（语义重建）。

\section{当前方法}

\subsection{核心想法}

现在的分支来自一个简单的设计理念：不要把同一种信号压进整个 SID 空间，而是让三层各自使用不同来源的图——L1 基于语义近邻路由保持粗分区稳定，L2 用多跳协同邻域做排序，L3 用局部转移精修补细粒度区分。

具体来说，当前主线做了三层约束设计：

\begin{itemize}
  \item L1：保持语义近邻路由，使第一层 SID 形成稳定、可学习的 coarse partition（粗分区）。L1 使用语义近邻图（semantic kNN graph），不引入协同信号，避免粗分区被协同图改写。
  \item L2：对齐多跳协同邻域，使第二层 SID 捕捉直接共现之外的短程协同结构。L2 使用 local multihop graph（局部多跳图）作为正样本来源，使用"语义近但图弱连"的物品对作为负样本。
  \item L3：对齐局部转移精修，使第三层 SID 保留同一物品族内部的细粒度区分。L3 与 L2 的图同源，但直接使用净化后的转移图，不做平方展开，保留最直接的短期共现关系。
\end{itemize}

\subsection{多粒度协同约束}

设 $A_{\mathrm{local}}$ 为从训练序列中构建的局部协同图：对每个训练样本，统计历史中每个物品到目标物品的转移频率来构建边关系，按时间近度加权（距离目标越近的历史物品权重越高），然后经过去噪、流行度修正和行归一化，最后每行保留 top-32 个邻居。我们将直接邻居与短程多跳邻居合并，形成中层使用的多跳协同图：

\[
A_{\mathrm{LMH}}
=
\row\left(A_{\mathrm{local}} + \alpha A_{\mathrm{local}}^2\right).
\]

直观上，$A_{\mathrm{local}}$ 捕捉直接共现/转移关系；$A_{\mathrm{local}}^2$ 是 $A_{\mathrm{local}}$ 与自身的矩阵乘积，表示经过一个中间物品的 2 跳间接关系，捕捉短程协同邻域。因此 $A_{\mathrm{LMH}}$ 同时保留了直接邻居和近距离间接邻居的结构。该图作为 L2 的正样本来源；L1 使用纯语义近邻图，L3 使用 local transition graph（局部转移图）。三层约束各有独立的图来源和损失形式，共同服务于同一个目标：让 SID 层级结构与协同层级结构对齐。

\subsection{层级 stop-gradient}

为了避免高层协同损失破坏低层路由，我们使用 previous-level stop-gradient（前层停梯度）：

\[
h_i^{(1)} = q_i^{(1)},
\]

\[
h_i^{(2)} = \sg(q_i^{(1)}) + q_i^{(2)},
\]

\[
h_i^{(3)} = \sg(q_i^{(1)}+q_i^{(2)}) + q_i^{(3)}.
\]

这样，每一层的损失主要作用在对应的 residual level（残差层级）上：L1 的语义 pull 不干扰 L2 的协同排序，L2 的 InfoNCE 不改写 L1 的粗分区。

\subsection{当前 tokenizer 目标函数}

当前主线的总损失为：

\[
\mathcal{L}
=
\mathcal{L}_{\mathrm{RQ}}
+\lambda_1\mathcal{L}_{\mathrm{L1}}
+\lambda_2\mathcal{L}_{\mathrm{L2}}
+\lambda_3\mathcal{L}_{\mathrm{L3}}.
\]

其中每个损失项使用不同的图来源和约束形式：

\begin{itemize}
  \item \textbf{L1 语义路由}（semantic routing）：在 L1 表示 $h_i^{(1)}=q_i^{(1)}$ 上施加 semantic graph（语义近邻图）的 pairwise pull 损失。拉近语义上相近的物品，形成稳定的粗分区：
  \[
  \mathcal{L}_{\mathrm{L1}}
  =\frac{1}{|\mathcal{B}|}\sum_{i,j} A^{\mathrm{sem}}_{ij}\,
  \lVert h_i^{(1)}-h_j^{(1)}\rVert_2^2.
  \]

  \item \textbf{L2 协同排序}（collaborative ranking）：在 L2 表示 $h_i^{(2)}=\sg(q_i^{(1)})+q_i^{(2)}$ 上施加 weighted graph InfoNCE 损失。正样本来自 local multihop graph，负样本来自预计算的"语义近但图弱连"物品对：
  \[
  \mathcal{L}_{\mathrm{L2},i}
  =-\log
  \frac{
  \sum_{j} A^{\mathrm{LMH}}_{ij}\exp(\cos(h_i^{(2)},h_j^{(2)})/\tau)
  }{
  \sum_{j} A^{\mathrm{LMH}}_{ij}\exp(\cos(h_i^{(2)},h_j^{(2)})/\tau)
  +\sum_{k} P^{\mathrm{weak}}_{ik}\exp(\cos(h_i^{(2)},h_k^{(2)})/\tau)
  }.
  \]

  \item \textbf{L3 局部精修}（local refinement）：在 L3 表示 $h_i^{(3)}=\sg(q_i^{(1)}+q_i^{(2)})+q_i^{(3)}$ 上施加 local transition graph（局部转移图）的 pairwise pull 损失。拉近在短期序列中直接共现的物品，提供细粒度区分：
  \[
  \mathcal{L}_{\mathrm{L3}}
  =\frac{1}{|\mathcal{B}|}\sum_{i,j} A^{\mathrm{local}}_{ij}\,
  \lVert h_i^{(3)}-h_j^{(3)}\rVert_2^2.
  \]
\end{itemize}

当前实现中，$\lambda_1=0.03$, $\lambda_2=0.01$, $\lambda_3=0.02$，温度 $\tau=0.1$。L2 的负样本对 $P^{\mathrm{weak}}$ 在训练前离线预计算。

\section{当前主要实验结果}

\subsection{下游 SFT 结果：与 baseline 的完整对比}

当前主线在 SFT/evaluate（监督微调/评测）上的结果如下。

\begin{center}
\begin{tabular}{lrrrr}
\toprule
Metric（指标） & @1 & @3 & @5 & @10 \\
\midrule
基线 NDCG & 0.067064 & 0.085008 & 0.093153 & 0.103720 \\
\textbf{本文方法 NDCG} & \textbf{0.070593} & \textbf{0.088131} & \textbf{0.094889} & \textbf{0.104383} \\
$\Delta$ NDCG & $+0.003530$ & $+0.003123$ & $+0.001736$ & $+0.000663$ \\
\midrule
基线 HR & 0.067064 & 0.098390 & 0.118244 & 0.150893 \\
\textbf{本文方法 HR} & \textbf{0.070593} & \textbf{0.100816} & 0.117362 & 0.146923 \\
$\Delta$ HR & $+0.003530$ & $+0.002427$ & $-0.000882$ & $-0.003971$ \\
\bottomrule
\end{tabular}
\end{center}

主要结论：

\begin{itemize}
  \item 相比 baseline，当前方法提升了 $\mathrm{NDCG@1/3/5/10}$。
  \item HR（命中率）在 @1 和 @3 提升，但 @5 和 @10 低于基线，说明当前方法更明显改善前排 排序质量，覆盖面仍有优化空间。
  \item 因此当前结果的核心价值是：tokenizer-side（分词器侧）协同层级注入已经带来稳定的 NDCG 提升。
\end{itemize}

\subsection{Tokenizer 结构观察}

我们观察到，成功点并不是 tokenizer collision（分词器碰撞）最低的点。更重要的是 SID 空间是否形成合理的层级结构。

\begin{center}
\small
\begin{tabular}{lccccc}
\toprule
\shortstack{Tokenizer\\（分词器）}
& \shortstack{活跃\\L1}
& \shortstack{唯一\\L1-L2}
& \shortstack{Collision\\（碰撞）}
& \shortstack{Top-5\\L1 cover}
& \shortstack{SFT\\NDCG@10} \\
\midrule
baseline & 48 & 2295 & 16 & 833 & 0.103720 \\
\textbf{本文方法} & \textbf{60} & \textbf{2330} & \textbf{16} & \textbf{707} & \textbf{0.104383} \\
\bottomrule
\end{tabular}
\end{center}

表中各列的含义：活跃 L1 指训练后被实际使用的第一层码本数量（共 256 个）；唯一 L1-L2 指实际出现的第一层-第二层前缀组合数，反映码本的分支多样性；Collision 指生成码本中发生冲突的物品数（共 3686 个物品）；Top-5 L1 cover 指最大的 5 个 L1 码所覆盖的物品总数，衡量 L1 的头部集中度。

从这个表可以看到：当前成功点的活跃 L1 为 60，接近原版 48；Top-5 L1 cover 为 707，低于原版 833，说明 L1 分布比原版更均匀一些。两者都表明较紧凑的 L1 可能更适合下游学习。

\subsection{错误分析}

对输出进行逐样本分析后，我们发现当前方法的收益更偏向排序质量，而不是命中覆盖。

与 baseline SFT 对比：

\[
\mathrm{Hit@10}_{\mathrm{new}} = 666,\qquad
\mathrm{Hit@10}_{\mathrm{orig}} = 684.
\]

虽然新方法命中样本更少，但它把更多正确目标排到了更靠前位置：

\[
\Delta \mathrm{NDCG@10}=+0.000663.
\]

这说明第二层的多跳协同结构已经被下游模型利用；同时，第一层粗路由和第三层转移精修仍有继续优化空间：

\begin{itemize}
  \item L1 无历史参照（目标物品的第一层码在用户历史中从未出现过）的样本收益很小；
  \item 高相似度物品仍然容易混淆，例如 3D 打印耗材的颜色、材料、品牌变体。
\end{itemize}

\end{document}
